\section*{第 \ref{ch:b1:systems-of-points} 章——质点组；刚体初步}

\begin{solution}{exo:b1:systems-of-points:1}
$x_G = (2.0 \times 0 + 3.0 \times 1.0)/5.0 = \qty{0.60}{m}$。地月系统：
$384000/82 = \qty{4700}{km}$，即距地心 \qty{4700}{km}，位于地球内部
（\qty{6370}{km}）：地球绕一点摆动，该点位于
其表面以下 \qty{1700}{km} 处。
\end{solution}

\begin{solution}{exo:b1:systems-of-points:2}
$4.0v = 0.010 \times 800$：$v = \qty{2.0}{m/s}$。子弹 $\tfrac12 \times 0.010
\times 800^2 = \qty{3.2}{kJ}$，步枪 \qty{8}{J}：来自火药的
化学能；轻的物体几乎获得全部能量（$E_k = p^2/2m$
，动量 $p$ 相同）。
\end{solution}

\begin{solution}{exo:b1:systems-of-points:3}
$v = 20000/2500 = \qty{8.0}{m/s}$；$E_k$：碰撞前 \qty{200}{kJ}，$\tfrac12 \times
2500 \times 64 = \qty{80}{kJ}$（碰撞后），损失 $60\%$。
\end{solution}

\begin{solution}{exo:b1:systems-of-points:4}
端点：$\int_0^L x^2(M/L)\dd x = ML^2/3$；中心：$\int_{-L/2}^{L/2}x^2(M/L)
\dd x = ML^2/12$；由惠更斯定理，取 $d = L/2$：$ML^2/12 + ML^2/4 = ML^2/3$。
\end{solution}

\begin{solution}{exo:b1:systems-of-points:5}
$E_k = \tfrac12 Mv^2(1 + k)$：圆柱体为 $1.5$，圆环为 $2$，均乘以 $\tfrac12 Mv^2$。
速率相同时，圆环的能量更多：在平地上它滚动得更远，
摩擦力相同，并能爬得更高（$h = v^2(1 + k)/2g$）。
\end{solution}

\begin{solution}{exo:b1:systems-of-points:6}
HCl：$\mu = 35m_H/36 = 0.972m_H$；原子：$\mu = m_e/(1 + 1/1836) =
0.99946m_e$。频率 $\propto 1/\sqrt\mu$：HCl 增加 $+1.4\%$，
氢增加 $+0.027\%$（H 与 D 光谱之间的同位素位移即由此而来）。
\end{solution}

\begin{solution}{exo:b1:systems-of-points:7}
$J = ML^2/3$，$d = L/2$：$T = 2\pi\sqrt{(ML^2/3)/(MgL/2)} = 2\pi\sqrt{2L/3g}
= \qty{1.64}{s}$；等效摆长 $2L/3 = \qty{0.67}{m}$。在
打击中心施加冲击，细杆将绕支点转动，支点处无冲击性反作用力；
击球时若击中球棒的“甜区”，
双手便不会感到震痛。
\end{solution}

\begin{solution}{exo:b1:systems-of-points:8}
$J = \tfrac12 MR^2 = \qty{1.0e-3}{kg.m^2}$；$C = 4\pi^2J/T^2 = \qty{9.9e-3}{N.m/rad}$。
\end{solution}

\begin{solution}{exo:b1:systems-of-points:9}
$\omega = \qty{314}{rad/s}$，$E = \tfrac12 J\omega^2 = \qty{4.9e5}{J}$；
$\Gamma = J\omega/t = \qty{314}{N.m}$；转数：$\tfrac12\omega t/2\pi = 250$。
\end{solution}

\begin{solution}{exo:b1:systems-of-points:10}
$a = g\sin\alpha/(1 + k)$：球体（$k = 2/5$）为 $0.714$，圆柱体（$1/2$）为
$0.667$，圆环（$1$）为 $0.500$，均乘以 $g\sin\alpha$：球体最先到达，圆环最后；
其中不含质量或半径。
\end{solution}

\begin{solution}{exo:b1:systems-of-points:11}
动量：$0.010 \times 400 = 2.01\,v$，$v = \qty{2.0}{m/s}$；再由能量：
$h = v^2/2g = \qty{0.20}{m}$。$E_k$：碰撞前 \qty{800}{J}，$\tfrac12 \times
2.01 \times 4.0 = \qty{4.0}{J}$（碰撞后），$99.5\%$ 转化为形变能与
热量。碰撞是非弹性的——能量在此过程中不守恒，
动量守恒（绳的张力竖直，且碰撞时间极短）。
\end{solution}

\begin{solution}{exo:b1:systems-of-points:12}
$m_1v_1 = m_1v_1' + m_2v_2'$ 与 $m_1v_1^2 = m_1v_1'^2 + m_2v_2'^2$ 联立，
得 $v_1 + v_1' = v_2'$（将能量方程除以动量方程），
由此得出各公式。传递给靶的能量比 $4m_1m_2/(m_1 + m_2)^2$：
质子 $100\%$；碳 $4 \times 12/169 = 28\%$；铅 $4 \times 207/208^2 =
1.9\%$。轻原子核经几次碰撞即可带走中子能量，
（水、石墨）；铅则需数百次。
\end{solution}

\begin{solution}{pb:b1:systems-of-points:1}
\textbf{1.} 在轴上，位于圆盘中心处，由对称性可得。

\textbf{2.} 柯尼希定理：$E_k = \tfrac12 Mv_G^2 + E_k^*$，在
质心参考系中，悠悠球绕其轴转动：$E_k^* = \tfrac12 J\omega^2$。

\textbf{3.} 半径为 $\rho$ 的圆环，质量为 $2M\rho\,\dd\rho/R^2$：$J = \int_0^R
\rho^2\cdot 2M\rho\,\dd\rho/R^2 = MR^2/2 = 0.5 \times 0.050 \times 9\times10^{-4}
= \qty{2.25e-5}{kg.m^2}$。

\textbf{4.} 细绳固定，并从轴上退绕：轴上与之接触的点
瞬时静止，故 $v_G = r\omega$。

\textbf{5.} $E_k = \tfrac12 Mv_G^2 + \tfrac12(\tfrac12 MR^2)(v_G/r)^2 = \tfrac12
Mv_G^2(1 + R^2/2r^2)$；$R^2/2r^2 = 9/(2 \times 0.16) = 28$：因子为 $29$——
几乎全部能量都在自旋中。

\textbf{6.} 重力 $Mg$ 向下，张力 $T$ 向上：$Ma_G = Mg - T$。

\textbf{7.} 只有张力对轴有力矩（力臂为 $r$）：
$J\dot\omega = Tr$。

\textbf{8.} $\dot\omega = a_G/r$：$T = Ja_G/r^2 = (R^2/2r^2)Ma_G$；于是
$Ma_G = Mg - (R^2/2r^2)Ma_G$，$a_G = g/29 = \qty{0.34}{m/s^2}$。

\textbf{9.} $T = M(g - a_G) = 0.050 \times 9.47 = \qty{0.47}{N}$，为
重力的 $97\%$：细绳几乎承受了全部重量。

\textbf{10.} $t = \sqrt{2h/a_G} = \qty{2.4}{s}$；$v_G = a_Gt = \qty{0.82}{m/s}$。
一颗石子：\qty{0.45}{s} 与 \qty{4.4}{m/s}。

\textbf{11.} $\omega = v_G/r = \qty{205}{rad/s} = \qty{1960}{rpm}$。

\textbf{12.} $\tfrac12 Mv_G^2 = \qty{0.017}{J}$，对比 $\tfrac12 J\omega^2 =
\qty{0.47}{J}$：自旋中的能量是下落时的 $28$ 倍。

\textbf{13.} $Mgh = \qty{0.49}{J}$；$\tfrac12 Mv_G^2 \times 29 = 0.5 \times 0.050
\times 0.67 \times 29 = \qty{0.49}{J}$。张力作用于
细绳上静止的那一点（细绳不运动）：功率为零。

\textbf{14.} 制动力沿细绳、经过轴线作用：
对该轴不产生力矩；角动量 $J\omega$ 保持不变。只有
切向力（轴上的摩擦力）才能减慢自旋。

\textbf{15.} $T \approx Mv_G/\Delta t + Mg = 0.050 \times 0.82/0.005 + 0.49 =
\qty{8.7}{N}$，为重力的十八倍——这正是
手指感到冲击的原因。

\textbf{16.} $J\dot\omega = -\Gamma$：$t = J\omega/\Gamma = 2.25\times10^{-5} \times
205/10^{-4} = \qty{46}{s}$。

\textbf{17.} 全部自旋能量 $\tfrac12 J\omega^2 = \qty{0.47}{J}$ 转化
为 $Mgh'$：$h' = \qty{0.96}{m}$，接近全部高度；实际上更少，
因为“睡眠”期间轴上的摩擦以及
冲击时的损耗。

\textbf{18.} $a_G$ 向下（悠悠球在上升过程中减速），而
$v_G$ 向上：$T = M(g - a_G) < Mg$——小于重力。

\textbf{19.} 当 $r = R$ 时：因子 $1 + 1/2 = 1.5$，$a_G = 2g/3 = \qty{6.5}{m/s^2}$
——下落快而自旋少：细轴才使其成为悠悠球。

\textbf{20.} 纯滚动：位于 $C$ 处的质点的
速度为零；一点静止的刚体，其速度场
是绕该点的转动（转轴过 $C$，角速度 $\omega$ 相同）。

\textbf{21.} $F$ 作用于轴的底部，高出桌面 $R - r$，
水平指向手，力矩 $F(R - r)$，以 $C$ 为参考，
转向使线轴朝手的方向滚动；重力与
支持力均过 $C$（无力矩）。

\textbf{22.} $(J + MR^2)\dot\omega = F(R - r)$，$a_G = R\dot\omega = FR(R - r)/
(\tfrac32 MR^2) = 2F(R - r)/3MR > 0$：指向手。

\textbf{23.} 当 $F$ 的作用线过 $C$ 时，即细线
从轴底部引出并与桌面交于 $C$：$\cos\beta = r/R$，
$\beta = 82^\circ$。角度更陡时，对 $C$ 的力矩反向，线轴向
远离手的方向滚动（$\beta$ 足够大时还会抬起）。

\textbf{24.} $a_G = 2 \times 0.20 \times 0.026/(3 \times 0.050 \times 0.030) =
\qty{2.3}{m/s^2}$；由 $Ma_G = F - f$ 得摩擦力（摩擦力与 $F$ 方向相反）：
$f = 0.20 - 0.050 \times 2.3 = \qty{0.085}{N}$；它须小于 $\mu_sMg
\approx \qty{0.2}{N}$（取 $\mu_s = 0.4$），即 $F \lesssim \qty{0.5}{N}$。

\textbf{25.} 质心：$G$ 的平动与张力；
绕轴的角动量：自旋与滚动方向；
能量：底部的速率与上升的高度。
\end{solution}
